% SEGMENT OLP-0032-S01
% Part:sets-functions-relations
% Chapter: size-of-sets
% Section: pairing-alt

\documentclass[../../../include/open-logic-section]{subfiles}

\begin{document}

\olfileid{sfr}{siz}{pai-alt}

\olsection{ஒரு மாற்று ஜோடியாக்கச் சார்பு}

\begin{explain}
$\Nat^2$ இன் வேறு சில பட்டியலாக்கங்களின் நேர்மாறுகளைக் கண்டறிவது
மேலும் எளிதாக இருக்கும். அவற்றில் ஒன்று இதோ.
பட்டியலாக்கத்தை ஓர் அட்டவணையில் காட்சிப்படுத்துவதற்குப் பதிலாக,
(தொடக்கத்தில்) வெற்றாக உள்ள இடங்களுடன் இணைக்கப்பட்ட நேர்ம முழுக்களின்
பட்டியலிலிருந்து தொடங்குக.
இந்த இடங்களை $\tuple{n,m}$ என்ற ஜோடிகளால் பின்வருமாறு
அடுத்தடுத்து நிரப்புவதாகக் கற்பனை செய்க.
முதல் இடத்தில் $0$ ஐக் கொண்ட ஜோடிகளிலிருந்து
(அதாவது $\tuple{0,m}$ என்ற ஜோடிகளிலிருந்து) தொடங்குக.
முதலாவதை (அதாவது $\tuple{0,0}$ ஐ) முதல் வெற்று இடத்தில் வைக்கவும்;
பின்னர் ஒரு வெற்று இடத்தைத் தாண்டி,
% LOCAL SOURCE CORRECTION TA-SIZ-LOCAL-002; see translation/size-notes.tex.
இரண்டாவதை (அதாவது $\tuple{0,1}$ ஐ) அடுத்த வெற்று இடத்தில் வைக்கவும்;
மீண்டும் ஒன்றைத் தாண்டவும்; இவ்வாறே தொடர்க.
நமது பட்டியலாக்கத்தின் (முழுமையடையாத) தொடக்கம் இப்போது பின்வருமாறு:
\[\small
\begin{array}{@{}c c c c c c c c c c c@{}}
\mathbf 1 & \mathbf 2 & \mathbf 3 & \mathbf 4 & \mathbf 5 & \mathbf 6 & \mathbf 7 & \mathbf 8 & \mathbf 9 & \mathbf{10} & \dots \\ \\
\tuple{0,0} &  & \tuple{0,1} &  & \tuple{0,2} &  & \tuple{0,3} & & \tuple{0,4} &  & \dots \\
\end{array}
\]
இன்னும் வெற்றாக உள்ள இடங்களில் $\tuple{1,m}$ என்ற ஜோடிகளைக் கொண்டு
இதை மீண்டும் செய்க; ஒவ்வொரு மாற்று வெற்று இடத்தையும் மீண்டும் தாண்டுக:
\[\small
\begin{array}{@{}c c c c c c c c c c c@{}}
\mathbf 1 & \mathbf 2 & \mathbf 3 & \mathbf 4 & \mathbf 5 & \mathbf 6 & \mathbf 7 & \mathbf 8 & \mathbf 9 & \mathbf{10} & \dots \\ \\
\tuple{0,0} & \tuple{1,0} & \tuple{0,1} &  & \tuple{0,2} & \tuple{1,1} &
\tuple{0,3} & & \tuple{0,4} &  \tuple{1,2} & \dots \\
\end{array}
\]
% SOURCE CORRECTION OLSIZ-003; see translation/size-notes.tex.
$\tuple{2,m}$, $\tuple{3,m}$ ஆகிய ஜோடிகள் முதலானவற்றையும்
இதே முறையில் உள்ளிடுக.
இவ்வாறு நிறைவடைந்த நமது பட்டியலாக்கம் பின்வருமாறு தொடங்குகிறது:
\[\small
\begin{array}{@{}cc c c c c c c c c c@{}}
\mathbf 1 & \mathbf 2 & \mathbf 3 & \mathbf 4 & \mathbf 5 & \mathbf 6 & \mathbf 7 & \mathbf 8 & \mathbf 9 & \mathbf{10} & \dots \\ \\
\tuple{0,0} & \tuple{1,0} & \tuple{0,1} & \tuple{2,0}  & \tuple{0,2} &
\tuple{1,1} & \tuple{0,3} & \tuple{3,0}  & \tuple{0,4} &  \tuple{1,2} & \dots \\
\end{array}
\]
மேலுள்ள அட்டவணையின் அறைகளுக்கு இந்தப் பட்டியலாக்கத்தின்படி
எண்களிட்டால், நேர்த்தியான வளைவழிக் கோடு கிடைக்காது;
பின்வரும் அமைப்பு கிடைக்கும்:
\[
\begin{array}{ c | c | c | c | c | c | c | c }
& \mathbf 0 & \mathbf 1 & \mathbf 2 & \mathbf 3 & \mathbf 4 & \mathbf 5 & \dots \\
\hline
\mathbf 0 & 1 & 3 & 5 & 7 & 9 & 11 & \dots \\
\hline
\mathbf 1 & 2 & 6 & 10 & 14 & 18 & \dots & \dots \\
\hline
\mathbf 2 & 4 & 12 & 20 & 28 & \dots & \dots & \dots \\
\hline
\mathbf 3 & 8 & 24 & 40 & \dots & \dots & \dots & \dots \\
\hline
\mathbf 4 & 16 & 48 & \dots & \dots & \dots & \dots & \dots \\
\hline
\mathbf 5 & 32 & \dots & \dots & \dots & \dots & \dots & \dots \\
\hline
\vdots & \vdots & \vdots & \vdots & \vdots & \vdots & \vdots & \ddots\\
\end{array}
\]
கிடைவரிசை $0$ இல் உள்ள ஜோடிகள் நமது பட்டியலாக்கத்தின்
ஒற்றைப்படை எண்களுடைய இடங்களில் உள்ளன.
அதாவது $\tuple{0,m}$ என்ற ஜோடி $2m+1$ ஆவது இடத்தில் உள்ளது.
இரண்டாவது கிடைவரிசையில் உள்ள $\tuple{1,m}$ என்ற ஜோடிகள்,
ஓர் ஒற்றைப்படை எண்ணின் இருமடங்கு எண் கொண்ட இடங்களில்,
குறிப்பாக $2 \cdot (2m+1)$ என்ற இடங்களில், உள்ளன.
மூன்றாவது கிடைவரிசையில் உள்ள $\tuple{2,m}$ என்ற ஜோடிகள்,
ஓர் ஒற்றைப்படை எண்ணின் நான்கு மடங்கான
$4 \cdot (2m+1)$ என்ற இடங்களில் உள்ளன; இவ்வாறே தொடர்கிறது.
ஒவ்வொரு கிடைவரிசைக்கும் $(2m+1)$ ஐப் பெருக்கும்
$1$, $2$, $4$, $8$, \dots ஆகிய காரணிகள்,
சரியாக $2$ இன் அடுக்குகளாகும்:
$1= 2^0$, $2 = 2^1$, $4 = 2^2$, $8 = 2^3$, \dots\@
உண்மையில், பொருத்தமான அடுக்குக்குறி எப்போதும் அந்த ஜோடியின்
முதல் உறுப்பு ஆகும்.
ஆகவே $\tuple{n,m}$ என்ற ஜோடிக்குரிய காரணி $2^n$.
இது $2^n \cdot (2m+1)$ என்ற பொதுவாய்பாட்டைத் தருகிறது.
ஆனால் இது ஜோடிகளை \emph{நேர்ம} முழுக்களுக்கு அனுப்பும் ஓர் இணைப்பு;
அதாவது, $\tuple{0,0}$ இன் இடம் $1$.
$0$ ஆவது இடத்திலிருந்து தொடங்க வேண்டுமானால்,
முடிவிலிருந்து $1$ ஐக் கழிக்க வேண்டும்.
இதனால் பின்வருவதைப் பெறுகிறோம்:
\end{explain}

% SEGMENT OLP-0032-S02
\begin{ex}
பின்வருமாறு கொடுக்கப்படும் $h\colon \Nat^2 \to \Nat$ என்ற சார்பு:
\[
h(n,m) = 2^n (2m+1) - 1
\]
இயல் எண்களின் ஜோடிகளைக் கொண்ட $\Nat^2$ கணத்திற்கான
ஒரு ஜோடியாக்கச் சார்பாகும்.
\end{ex}

% SEGMENT OLP-0032-S03
\begin{explain}
அதற்கேற்ப, $\Nat^2$ இன் நமது இரண்டாவது பட்டியலாக்கத்தில்,
$\tuple{0,0}$ இன் குறியெண் $h(0,0) = 2^0(2\cdot 0+1) - 1 = 0$;
$\tuple{1,2}$ இன் குறியெண் $2^{1} \cdot (2 \cdot 2 + 1) - 1 = 2
\cdot 5 - 1 = 9$;
$\tuple{2,6}$ இன் குறியெண் $2^{2} \cdot (2
\cdot 6 + 1) - 1 = 51$.
\end{explain}

% SEGMENT OLP-0032-S04
சில சமயங்களில், குறியாக்கம் மேற்கோர்த்தல் பண்புடையதாக இருக்க வேண்டும்
என்று கோராமல், இயல் எண்களின் ஜோடிகள் $\Nat^2$ ஐக்
குறியாக்கினாலே போதுமானது.
இத்தகைய குறியாக்கங்களின் நேர்மாறுகள் பகுதிச் சார்புகளாக மட்டுமே இருக்கும்.

% SEGMENT OLP-0032-S05
\begin{ex}
பின்வருமாறு கொடுக்கப்படும் $j\colon \Nat^2 \to \Nat^+$ என்ற சார்பு:
\[
j(n,m) = 2^n3^m
\]
$\Nat^2 \to \Nat$ என்ற !!a{injective} சார்பாகும்.
\end{ex}

\end{document}
